% !TEX program = pdflatex
% ============================================================
% Guía de estudio — Interrogación 1 (ICS2122 Capstone, Grupo 29)
% Archivo autocontenido: se sube tal cual a Overleaf y compila con pdfLaTeX.
%
% Para practicar sin ver las respuestas modelo de la sección 7,
% cambia \respuestastrue por \respuestasfalse (línea de abajo).
% ============================================================
\documentclass[11pt,a4paper,parskip=half]{scrartcl}

\newif\ifrespuestas
\respuestastrue   % \respuestasfalse = deja espacio en blanco para responder a mano

% ---------- base ----------
\usepackage[utf8]{inputenc}
\usepackage[T1]{fontenc}
\usepackage{lmodern}
\usepackage[spanish,es-tabla,es-noshorthands,es-nodecimaldot]{babel}
\usepackage[a4paper,margin=2.3cm]{geometry}
\usepackage{microtype}
\usepackage{amsmath,amssymb}
\usepackage{enumitem}
\usepackage{booktabs,longtable,tabularx,array}

% ---------- paleta del proyecto (misma de 04-documentos/estilo.tex) ----------
\usepackage[table]{xcolor}
\definecolor{ink}{RGB}{0,62,92}
\definecolor{inklight}{RGB}{92,132,153}
\definecolor{rulegray}{RGB}{205,214,219}
\definecolor{critred}{RGB}{158,38,32}
\definecolor{amber}{RGB}{176,98,16}
\definecolor{slate}{RGB}{104,112,118}
\definecolor{softrow}{RGB}{242,246,248}

\usepackage{tikz}
\usetikzlibrary{arrows.meta,positioning,calc,fit,backgrounds}
\usepackage{tcolorbox}
\usepackage[colorlinks=true,linkcolor=ink,urlcolor=inklight]{hyperref}

% ---------- títulos ----------
\setkomafont{disposition}{\color{ink}\sffamily\bfseries}
\setkomafont{section}{\color{ink}\sffamily\Large\bfseries}
\setkomafont{subsection}{\color{ink}\sffamily\large\bfseries}
\setkomafont{subsubsection}{\color{ink}\sffamily\normalsize\bfseries}

% ---------- listas compactas ----------
\setlist{itemsep=2pt,topsep=3pt,parsep=0pt}

% ---------- tablas ----------
\newcolumntype{L}[1]{>{\raggedright\arraybackslash}p{#1}}
\newcolumntype{R}[1]{>{\raggedleft\arraybackslash}p{#1}}
\newcolumntype{Y}{>{\raggedright\arraybackslash}X}
% \leavevmode evita que \color desalinee la primera línea de una celda p{}
\newcommand{\hd}[1]{\leavevmode{\sffamily\bfseries\color{ink}#1}}
\renewcommand{\arraystretch}{1.3}
\setlength{\tabcolsep}{5pt}
\setlength{\LTpre}{6pt}
\setlength{\LTpost}{6pt}

% nombres de archivos y columnas del GTFS
\newcommand{\arch}[1]{\texttt{#1}}

% ---------- cajas ----------
\newtcolorbox{clave}{colback=softrow, colframe=ink, boxrule=0pt, leftrule=3pt,
  sharp corners, left=8pt, right=8pt, top=5pt, bottom=5pt, fontupper=\small}
\newtcolorbox{ojo}{colback=amber!7, colframe=amber, boxrule=0pt, leftrule=3pt,
  sharp corners, left=8pt, right=8pt, top=5pt, bottom=5pt, fontupper=\small}

% ---------- preguntas de práctica ----------
\newcounter{preg}
\newcommand{\pregunta}[1]{%
  \refstepcounter{preg}%
  \par\addvspace{10pt}%
  {\parfillskip=0pt\noindent\color{rulegray}\rule{\linewidth}{0.6pt}\par}\nopagebreak
  \noindent{\sffamily\bfseries\color{ink}P\thepreg.\ #1}\par\nopagebreak}
\newcommand{\respuesta}[1]{%
  \ifrespuestas
    {\small\noindent\textcolor{inklight}{\sffamily\bfseries Respuesta modelo.}\ #1\par}%
  \else
    \par\vspace{6cm}\penalty0 % permite cortar página después del espacio
  \fi}

\begin{document}

% ============================================================
\begin{center}
  {\sffamily\small\color{inklight}ICS2122 Taller de Investigación Operativa (Capstone) $\cdot$ Grupo 29}\\[6pt]
  {\sffamily\bfseries\LARGE\color{ink}Guía de estudio --- Interrogación 1}\\[6pt]
  {\sffamily\color{slate}Planificación de la operación de buses eléctricos en la Red Metropolitana de Movilidad}\\[4pt]
  {\sffamily\small\color{slate}Preparada el 28 de septiembre de 2026 $\cdot$ Prueba: martes 29 de septiembre, 9:40}
\end{center}

\vspace{4pt}
{\small\tableofcontents}

% ============================================================
\section*{Cómo es la prueba}
\addcontentsline{toc}{section}{Cómo es la prueba}

Martes 29 de septiembre, 9:40, en horario de clases. Dura cerca de una hora, es en lápiz y papel y trae 3 o 4 preguntas de desarrollo: unos 12 a 15 minutos por pregunta.

\begin{center}
\small
\begin{tabularx}{\linewidth}{@{}YY@{}}
\toprule
\hd{Entra} & \hd{No entra} \\
\midrule
El problema, los datos, los análisis y las reflexiones del informe y la presentación & El modelo matemático propio de la Entrega 2 \\
Justificar cualquier decisión que tomamos & Detalles de implementación que todavía no existen \\
Aplicar una metodología del informe a nuestro caso, o a una variante del problema & \\
\bottomrule
\end{tabularx}
\end{center}

Habrá preguntas generales (iguales para todos los grupos, pero se responden con nuestro proyecto) y preguntas específicas de buses eléctricos RED.

\textbf{Fuentes de esta guía:} el informe entregado (Grupo 29, 11 de septiembre) y la presentación (3 de septiembre). Cuando las dos difieren, \textbf{manda el informe}: es posterior y es el documento corregido. Las diferencias están en la sección~\ref{sec:decisiones}.

\begin{clave}
\textbf{Plantilla para cualquier respuesta de desarrollo} (cuatro movimientos, uno o dos párrafos cada uno):
\begin{enumerate}
  \item \textbf{Qué es:} defínelo en una frase.
  \item \textbf{En nuestro problema:} amárralo a una cifra concreta (64.502 corridas, 6.601 buses, 700 cupos, 225 km).
  \item \textbf{Por qué lo hicimos así:} la justificación, idealmente contra la alternativa descartada.
  \item \textbf{Qué sacrificamos:} la limitación o el supuesto, y cómo lo mediríamos o corregiríamos.
\end{enumerate}
\end{clave}

Plan sugerido: hoy lee las secciones del problema a las decisiones y responde a mano dos preguntas de práctica con reloj (en el preámbulo puedes esconder las respuestas modelo). Mañana temprano, solo el repaso final.

% ============================================================
\section{El problema en una página}
\label{sec:problema}

Hay que armar el plan diario de la flota eléctrica de la Red Metropolitana de Movilidad (RED): qué bus cubre cada corrida, en qué orden, y cuándo y dónde recarga. Todo para cubrir las 64.502 corridas de un día laboral al menor costo diario.

\textbf{Por qué no es el problema de siempre.} Programar un bus diésel es decidir una sola cosa: la secuencia de viajes. El combustible se carga en minutos, a cualquier hora, y nadie compite por la bomba. En un bus eléctrico la batería no alcanza para la jornada, recargar toma del orden de 105 minutos, solo se puede en un electroterminal y el enchufe tiene cupo.

\begin{clave}
\emph{El plan de viajes y el plan de carga se deciden juntos o no se deciden.} (Frase de la presentación.)
\end{clave}

\begin{minipage}[t]{0.48\linewidth}
\textbf{Qué se decide}
\begin{itemize}
  \item Cuántos buses usar (la flota es una decisión, no un dato).
  \item Qué corridas hace cada bus y en qué orden, incluidos los traslados en vacío entre una y otra.
  \item Cuándo, dónde y cuánto recarga cada bus.
\end{itemize}
\end{minipage}\hfill
\begin{minipage}[t]{0.48\linewidth}
\textbf{Qué lo restringe}
\begin{itemize}
  \item Cobertura: cada corrida se cubre exactamente una vez.
  \item Tiempo y espacio: un bus toma la corrida siguiente solo si alcanza a terminar la actual, trasladarse y llegar a la hora.
  \item Energía: el estado de carga (SOC) nunca baja del 10\,\%.
  \item Infraestructura: 5 electroterminales con 700 cupos en total.
\end{itemize}
\end{minipage}

\textbf{Objetivo.} Un plan factible al menor costo diario. Los costos entregados son 250 USD por bus usado, 0,5 USD por km, 0,10 a 0,22 USD por kWh según la hora, 0,03 USD por minuto de espera y 5 USD por actividad de carga.

\textbf{Entregable.} Una herramienta computacional reproducible que entregue la asignación de corridas y recargas, la flota usada y los costos.

\textbf{Alcance que fijamos}
\begin{itemize}
  \item Las frecuencias del GTFS son oferta exógena: no se eliminan corridas ni se cambian frecuencias para ahorrar.
  \item Flota irrestricta: los 1.200 buses de \arch{parameters.csv} son una referencia y no un tope, por aclaración del equipo docente.
  \item Laboral, sábado y domingo o festivo se estudian por separado. El caso principal es el día laboral.
\end{itemize}

\textbf{Cómo se llama en la literatura.} \emph{Electric Vehicle Scheduling Problem} (E-VSP): la extensión eléctrica del problema clásico de programación de flota (VSP), según la revisión de Perumal, Lusby y Larsen (2022) que citamos en la presentación.

% ============================================================
\section{Los datos}
\label{sec:datos}

Del GTFS de Santiago reconstruimos 64.502 corridas laborales de buses RED. De ellas sale una cota inferior de 6.601 buses, que es un piso y no una flota factible.

\begin{center}
\small
\begin{tabularx}{\linewidth}{@{}L{2.9cm}L{5.4cm}Y@{}}
\toprule
\hd{Familia} & \hd{Archivos} & \hd{Qué aporta} \\
\midrule
Programación & \arch{routes}, \arch{trips}, \arch{frequencies}, \arch{stop\_times} & Recorridos, patrones, ventanas, headways, duración \\
Geografía & \arch{stops}, \arch{shapes}, \arch{shape\_distances} & Paraderos, trazados, km de servicio \\
Bus eléctrico & \arch{vehicles}, \arch{parameters} & Batería, autonomía, potencia de carga, límites de SOC, costos \\
Electroterminales & \arch{depots}, \arch{charger\_capacity}, \arch{charging\_activities} & Ubicación, capacidad física y nominal de carga \\
Energía & \arch{electricity\_prices} & Precio por período horario \\
\bottomrule
\end{tabularx}
\end{center}

\subsection*{Paso 1: filtrar}
El GTFS trae buses, metro, trenes y aeropuerto. Filtramos \arch{agency\_id = RM} y \arch{route\_type = 3} (bus) y propagamos el filtro por identificadores a las demás tablas. Quedan 417 de 427 recorridos y 14.464 de 26.137 patrones (6.620 laborales, 4.215 de sábado y 3.629 de domingo o festivo).

\subsection*{Paso 2: expandir frecuencias}
Todas las filas RED de \arch{frequencies.txt} tienen \arch{exact\_times = 0}: el archivo da reglas de frecuencia, no horarios. Generamos las salidas así:
\[
  t_n = \text{start\_time} + n \cdot \text{headway}, \qquad t_n < \text{end\_time}.
\]
Ejemplo de la presentación: ventana 07:00--08:00 con headway de 900\,s da 4 salidas (07:00, 07:15, 07:30, 07:45). Son una reconstrucción representativa, no horarios observados.

\subsection*{Paso 3: duración, distancia y energía de cada corrida}
La duración es la llegada a la última parada menos la salida de la primera:
\[
  T_i = \text{llegada}_{i,\,\text{última}} - \text{salida}_{i,\,\text{primera}}.
\]
La distancia viene de \arch{shape\_distances.csv} por \arch{shape\_id} (839 trazados, todos con distancia). El consumo no venía en los datos, así que lo estimamos lineal:
\[
  e = \frac{350\ \text{kWh}}{250\ \text{km}} = 1{,}4\ \text{kWh/km}.
\]
Con SOC mínimo de 10\,\% quedan 315\,kWh utilizables, o sea unos 225\,km de autonomía útil.

\subsection*{Resultado: la escala de un día}
\begin{center}
\small
\begin{tabular}{@{}lrrrrrc@{}}
\toprule
\hd{Día} & \hd{Patrones} & \hd{Corridas} & \hd{Km de servicio} & \hd{Energía (MWh)} & \hd{Cota de flota} & \hd{Pico} \\
\midrule
Laboral & 6.620 & 64.502 & 1.345.413 & 1.883,6 & 6.601 & 08:00 \\
Sábado & 4.215 & 43.619 & 924.134 & 1.293,8 & 3.860 & 15:00 \\
Domingo o festivo & 3.629 & 37.377 & 794.792 & 1.112,7 & 2.900 & 18:00 \\
\bottomrule
\end{tabular}
\end{center}
Una corrida laboral promedio dura 78,6\,min, recorre 20,9\,km y gasta 29,2\,kWh.

\subsection*{La cota de flota}
Contamos cuántas corridas están activas a la vez en cada minuto. El máximo laboral es 6.601 a las 08:00, así que ningún plan puede usar menos de 6.601 buses:
\[
  F_L \geq 6{.}601, \qquad F_S \geq 3{.}860, \qquad F_D \geq 2{.}900.
\]
Es condición \textbf{necesaria, no suficiente}: ignora traslados en vacío, recargas y regulación, así que la flota real será mayor. Tampoco la energía agregada asegura nada: cada batería es de un bus, y una cadena puede necesitar recarga aunque al sistema completo le sobre energía.

\textbf{Cifras derivadas que usa la presentación}
\begin{itemize}
  \item 203,9\,km por bus y día: 1.345.413\,km repartidos en unos 6.600 buses.
  \item 91\,\%: esos 203,9\,km contra los 225\,km útiles. Cada bus necesita casi una recarga diaria (0,91).
  \item 105\,min por recarga completa: 315\,kWh a 180\,kW.
\end{itemize}

\textbf{Infraestructura.} Cinco electroterminales con 700 cupos físicos: Vespucio Norte (150), El Conquistador (180), Los Espinos (120), La Reina (100) y Santa Rosa (150). El bus de referencia tiene 350\,kWh, 250\,km nominales y carga a 180\,kW.

\subsection*{Observado, derivado o supuesto}
Pregunta típica: saber distinguirlos.
\begin{center}
\small
\begin{tabular}{@{}ll@{}}
\toprule
\hd{Información} & \hd{Tipo} \\
\midrule
417 rutas, 6.620 patrones laborales & Observado tras el filtro \\
64.502 corridas, 6.601 simultáneas & Derivado (expansión de frecuencias) \\
Batería de 350\,kWh & Parámetro entregado \\
Consumo de 1,4\,kWh/km & Derivado y supuesto (lineal) \\
SOC inicial de 100\,\% & Supuesto (no viene en los datos) \\
\bottomrule
\end{tabular}
\end{center}

\subsection*{Limitaciones que declaramos}
\begin{itemize}
  \item \arch{exact\_times = 0}: las salidas representan frecuencias, no horarios reales.
  \item Duraciones nominales, sin congestión, accidentes, fallas ni mantenciones.
  \item Sin tiempo de detención en paraderos: llegada y salida coinciden.
  \item Sin demanda de pasajeros: por eso las frecuencias se toman como exógenas.
  \item Buses no individualizados y sin SOC inicial.
  \item Traslados en vacío no entregados: al reconstruirlos geográficamente son datos derivados.
  \item Consumo lineal: ignora pendiente, climatización, carga de pasajeros, tráfico y degradación.
  \item No se puede validar la capacidad eléctrica real de los electroterminales.
\end{itemize}

\begin{clave}
La conclusión del informe, que conviene citar tal cual: estas limitaciones obligan a \textbf{distinguir factibilidad dentro del modelo de factibilidad operacional real}.
\end{clave}

% ============================================================
\section{Dificultad, objetivo y KPIs}
\label{sec:kpis}

El problema es difícil por tres razones a la vez: su tamaño, el acoplamiento entre viajes y energía, y la competencia por la carga. Por eso minimizamos costo con la flota como palanca principal, y medimos contra una cota.

\subsection*{Por qué es difícil (lámina 4 de la presentación)}
\begin{center}
\small
\begin{tabularx}{\linewidth}{@{}L{2.4cm}L{4.6cm}Y@{}}
\toprule
\hd{Golpe} & \hd{Cifra} & \hd{Qué significa} \\
\midrule
Tamaño & 1.777 millones de arcos de encadenamiento posibles & Solo el 23\,\% de los pares de corridas es geográficamente factible y apenas el 3,1\,\% es barato \\
Acoplamiento & 91\,\% de la autonomía (203,9 de 225\,km) & Casi una recarga por bus al día, y cada una le quita 105\,min a la operación \\
Competencia & 6.601 buses simultáneos a las 08:00 & Contra 700 cupos de carga en 5 electroterminales \\
\bottomrule
\end{tabularx}
\end{center}

Remate de la lámina: el espacio de soluciones es inmenso y casi todo es infactible. Encontrar un plan que simplemente funcione ya es difícil, y queremos el mejor.

\begin{ojo}
\textbf{Argumento teórico que suma puntos} (conocimiento general del área, no está en el informe): el VSP con un solo depósito se resuelve en tiempo polinomial como flujo de costo mínimo. Con varios depósitos, o con límites de autonomía, pasa a ser NP-difícil. Nosotros tenemos las dos cosas.
\end{ojo}

\subsection*{Objetivo: dónde está el costo}
Un piso de costo calculado desde los datos da 2.512.204\,USD/día:
\begin{center}
\small
\begin{tabular}{@{}lrr@{}}
\toprule
\hd{Componente} & \hd{USD/día} & \hd{Parte del piso} \\
\midrule
Flota (250\,USD $\times$ bus) & 1.650.750 & 66\,\% \\
Kilómetros (0,5\,USD/km) & 673.011 & 27\,\% \\
Energía a 0,10\,USD/kWh & 188.443 & 7\,\% \\
\midrule
Total & 2.512.204 & 100\,\% \\
\bottomrule
\end{tabular}
\end{center}
Si toda la energía se cargara a 0,22\,USD/kWh, ese ítem sube a 414.575. La conclusión de la lámina: \textbf{el objetivo primario es usar menos buses, no gastar menos energía.} Cada bus que se ahorra vale 250\,USD al día; la diferencia total entre cargar caro y barato es de unos 226.000\,USD al día.

\subsection*{KPIs del informe (tabla 1.5, versión final)}
\begin{center}
\small
\begin{tabularx}{\linewidth}{@{}L{3.2cm}L{5.4cm}Y@{}}
\toprule
\hd{Dimensión} & \hd{Indicador} & \hd{Qué mide} \\
\midrule
Flota & Número de buses utilizados & Vehículos que requiere el plan \\
Costo & Costo operacional diario total & Todos los componentes económicos \\
Eficiencia operacional & \% de km recorridos sin pasajeros & Peso de los reposicionamientos \\
Infraestructura & Máxima ocupación simultánea de carga por electroterminal & Dónde y cuándo se aprieta la carga \\
Energía & Costo energético diario & Efecto de cuántas recargas y a qué hora \\
\bottomrule
\end{tabularx}
\end{center}

Tres ideas que el informe deja claras y que hay que saber defender:
\begin{enumerate}
  \item \textbf{Cobertura del 100\,\% no es KPI}, es condición de factibilidad. Un plan que deja una corrida sin cubrir no es ``un poco peor'': no es un plan.
  \item \textbf{No fijamos metas arbitrarias.} Sin una solución base factible no hay referencia comparable. Las metas saldrán de una solución base sencilla; las cotas inferiores solo dan contexto.
  \item \textbf{La brecha contra la cota no es brecha de optimalidad.} Para la flota:
  \[
    \text{Gap}_F = \frac{F_{sol} - F_{LB}}{F_{LB}}.
  \]
  Mide la distancia al límite teórico de 6.601. No es un gap de optimalidad mientras no se demuestre que la cota es ajustada, y no lo es, porque ignora traslados y recargas.
\end{enumerate}

\textbf{Los KPIs de la presentación eran otros.} En la lámina ``Con qué vamos a medir'' había valores de referencia: flota 6.601; ocupación de terminales 90\,\% (630 de 700 cupos); costo total 2.545.946\,USD; porcentaje del tiempo de un bus fuera de servicio bajo 20\,\%. En el informe se quitaron las metas por el punto 2 de arriba. Si te preguntan por qué cambiaron, esa es la respuesta.

% ============================================================
\section{Metodologías}
\label{sec:metodologias}

Comparamos tres metodologías y elegimos la red tiempo-espacio con descomposición de flujo en tres fases: sacrifica la garantía de optimalidad a cambio de una escala demostrada (10.710 viajes) cercana a la nuestra.

\begin{center}
\footnotesize
\begin{tabularx}{\linewidth}{@{}L{2.0cm}YYY@{}}
\toprule
 & \hd{MILP (modelo único)} & \hd{Generación de columnas} & \hd{Red tiempo-espacio + descomposición} \\
\midrule
\hd{Idea} & Una formulación con todo: cobertura, continuidad, flota, batería, capacidad. Se resuelve con \emph{branch and bound} & Cada columna es la jornada completa de un bus. Maestro de cobertura + subproblema de \emph{pricing} & Flujo agregado sin batería, luego se parte en rotaciones y se validan con energía \\
\hd{Qué garantiza} & Óptimo global & Cota inferior de calidad (relajación lineal del maestro) & Siempre una solución factible; sin garantía de óptimo \\
\hd{Escala publicada} & No cierra ni instancias de 150 viajes & 816 viajes (de Vos et al., 2024) & 10.710 viajes (Olsen et al., 2020) \\
\hd{Por qué sí o no} & Cerca de $5\times10^{8}$ variables binarias al expandir por SOC y electroterminal. Si no termina, no entrega nada útil & Nuestra instancia laboral es 79 veces más grande; el cuello de botella es el \emph{pricing}, que corre cientos o miles de veces & \textbf{Elegida.} Con 5 sectores de unos 12.000 viajes quedamos en la escala comprobada \\
\hd{Fuente} & Alvo, Angulo y Klapp (2021) & de Vos, van Lieshout y Dollevoet (2024) & Olsen, Kliewer y Wolbeck (2020) \\
\bottomrule
\end{tabularx}
\end{center}

\subsection{Generación de columnas, bien explicada}
La preguntan mucho aunque no la elegimos.
\begin{itemize}
  \item \emph{Maestro}: un problema de cobertura de conjuntos. Solo quedan dos familias de restricciones: cada corrida cubierta una vez, y la capacidad de carga por electroterminal y bloque horario.
  \item \emph{Pricing}: un camino más corto con restricciones de recursos sobre una red expandida por niveles de SOC. Cada nodo dice dónde y cuándo está el bus y cuánta energía le queda. Busca la jornada de costo reducido más negativo.
  \item Si encuentra una con costo reducido negativo, la agrega y repite. Si no hay ninguna, la relajación lineal está resuelta al óptimo, aunque solo se haya generado una fracción ínfima de columnas.
  \item Por qué la cota es buena: cada columna ya es factible por construcción, así que la relajación solo combina jornadas ejecutables.
  \item La solución suele salir fraccionaria: hay que usar heurísticas de integralidad o \emph{branch-and-price}.
\end{itemize}

\subsection{La metodología escogida}
Se apoya en Olsen et al. (2020). Su idea central son las \textbf{líneas de tiempo}: los arcos representan el paso del tiempo en cada punto, así que no hace falta un arco directo entre dos corridas si la misma conexión se logra esperando y trasladando por tramos. Eso elimina hasta el 97\,\% de los arcos de traslado en vacío, clave frente a los 1.777 millones de encadenamientos posibles.

\textbf{Dos descomposiciones antes de resolver}
\begin{itemize}
  \item \emph{Espacial}: las 417 rutas se reparten en 5 subproblemas, uno por electroterminal, según proximidad. Cada uno queda en unos 12.000 viajes. Es coherente con el MDVSP clásico (Kliewer, Mellouli y Suhl, 2006), que arma una red por depósito porque el acoplamiento entre depósitos es débil cuando cada uno tiene su propia capacidad.
  \item \emph{Temporal}: tres ejecuciones encadenadas. El día laboral debe terminar a las 23:59 en un estado igual o mejor que el de las 00:00, así se replica 5 veces sin volver a resolverse. El sábado hereda el estado final del laboral. El domingo hereda el del sábado y debe cerrar en el estado que el lunes supuso. El domingo, con menos demanda, da margen para recargar y reposicionar.
\end{itemize}

\begin{figure}[htbp]
\centering
\begin{tikzpicture}[
  font=\sffamily\small,
  caja/.style={draw=slate!80, rounded corners=3pt, fill=white, inner sep=5pt, line width=0.6pt},
  sector/.style={caja, text width=2.45cm, align=center, anchor=north},
  fase/.style={caja, text width=4.45cm, align=left, anchor=north},
  banda/.style={draw=rulegray, rounded corners=4pt, fill=softrow, inner sep=7pt},
  titulob/.style={font=\sffamily\bfseries\small, text=ink, anchor=north west, inner sep=0pt},
  flecha/.style={-{Stealth[length=2.4mm,width=1.8mm]}, draw=slate, line width=0.7pt},
  nota/.style={font=\sffamily\footnotesize, text=slate},
]
% entrada
\node[caja, text width=7.6cm, align=center] (in) at (0,0)
  {64.502 corridas de un día laboral $\cdot$ 417 rutas};

% descomposición espacial
\node[titulob] (tE) at (-7.55,-1.2)
  {Espacial: una red por electroterminal, unos 12.000 viajes cada una};
\node[sector] (s1) at (-6.2,-1.8) {Vespucio Norte\\{\footnotesize\color{slate}150 cupos}};
\node[sector] (s2) at (-3.1,-1.8) {El Conquistador\\{\footnotesize\color{slate}180 cupos}};
\node[sector] (s3) at (0,-1.8)    {Los Espinos\\{\footnotesize\color{slate}120 cupos}};
\node[sector] (s4) at (3.1,-1.8)  {La Reina\\{\footnotesize\color{slate}100 cupos}};
\node[sector] (s5) at (6.2,-1.8)  {Santa Rosa\\{\footnotesize\color{slate}150 cupos}};
\begin{scope}[on background layer]
  \node[banda, fit=(tE)(s1)(s5)] (bE) {};
\end{scope}

% tres fases
\node[titulob] (tF) at (-7.55,-4.1) {En cada sector y cada día: tres fases};
\node[fase] (f1) at (-5.25,-4.7)
  {\textbf{Fase 1 $\cdot$ Flujo relajado}\\
   {\color{slate}Red tiempo-espacio agregada}\\
   {\color{slate}Cubre cada viaje una vez}\\
   \textbf{Relaja batería y carga}\\
   {\color{slate}Entrega flota y km en vacío}};
\node[fase] (f2) at (0,-4.7)
  {\textbf{Fase 2 $\cdot$ Rotaciones}\\
   {\color{slate}Parte el flujo en jornadas}\\
   {\color{slate}Electroterminal: maximiza}\\
   {\color{slate}la espera mínima}\\
   {\color{slate}Terminal: equilibra consumo}};
\node[fase] (f3) at (5.25,-4.7)
  {\textbf{Fase 3 $\cdot$ Validación SOC}\\
   {\color{slate}Sigue el SOC tramo a tramo}\\
   {\color{slate}Inserta cargas en holguras}\\
   \textbf{Bajo 10\,\%: corta la}\\
   \textbf{rotación y agrega un bus}};
\coordinate (fL) at (bE.west |- f1.south);
\coordinate (fR) at (bE.east |- f1.south);
\begin{scope}[on background layer]
  \node[banda, inner xsep=0pt, fit=(tF)(f1)(f3)(fL)(fR)] (bF) {};
\end{scope}

% plan
\node[caja, draw=ink, line width=1.1pt, fill=ink!7, text width=8.8cm, align=center, anchor=north]
  (plan) at (0,-8.25)
  {\textbf{Plan factible: bus por corrida, recargas, flota y costo}\\
   {\footnotesize KPIs contra la cota de 6.601 buses y una solución base}};

% descomposición temporal
\node[titulob] (tT) at (-7.55,-10.45) {Temporal: tres ejecuciones encadenadas por semana};
\node[fase] (d1) at (-5.25,-11.05)
  {\textbf{Laboral}\\
   {\color{slate}se replica 5 veces}\\
   {\color{slate}estado 23:59 $\geq$ 00:00}};
\node[fase] (d2) at (0,-11.05)
  {\textbf{Sábado}\\
   {\color{slate}hereda el estado final}\\
   {\color{slate}del día laboral}};
\node[fase] (d3) at (5.25,-11.05)
  {\textbf{Domingo}\\
   {\color{slate}hereda el del sábado}\\
   {\color{slate}más holgura para recargar}};
\path (d3.south) ++(0,-0.45) coordinate (loopY);
\node[nota, anchor=north] (lc) at (0,0 |- loopY) [yshift=-3pt]
  {el domingo debe cerrar en el estado que el lunes supuso};
\coordinate (tL) at (bE.west |- lc.south);
\coordinate (tR) at (bE.east |- lc.south);
\begin{scope}[on background layer]
  \node[banda, inner xsep=0pt, fit=(tT)(d1)(d3)(lc)(tL)(tR)] (bT) {};
\end{scope}

% conectores
\draw[flecha] (in.south) -- (in.south |- bE.north);
\draw[flecha] (bE.south) -- (bE.south |- bF.north);
\draw[flecha] (f1.east) -- (f2.west);
\draw[flecha] (f2.east) -- (f3.west);
\draw[flecha] (bF.south) -- (bF.south |- plan.north);
\draw[flecha] (plan.south) -- node[right, nota, xshift=2pt]
  {el estado final pasa a la ejecución siguiente} (plan.south |- bT.north);
\draw[flecha] (d1.east) -- (d2.west);
\draw[flecha] (d2.east) -- (d3.west);
\draw[flecha] (d3.south) -- (loopY) -| (d1.south);
\end{tikzpicture}
\caption{La metodología escogida. Las dos descomposiciones achican el problema antes de resolverlo. Dentro de cada sector, la Fase~1 ignora la batería y la Fase~3 la recupera, a costa de la optimalidad.}
\label{fig:metodologia}
\end{figure}

\subsection{Las tres fases}

\textbf{Fase 1: flujo relajado.} Nodos: eventos de llegada y salida. Arcos: viajes de servicio, traslados en vacío, esperas y entradas y salidas de electroterminales. Se minimiza el costo de vehículos con conservación de flujo y cada viaje cubierto una vez. \textbf{Se relajan la batería y la capacidad de carga.} Entrega el flujo por arco, la flota mínima del sector sin energía, los km en vacío y los buses inactivos por lugar y hora.

\textbf{Fase 2: descomposición en rotaciones.} El flujo dice cuántos buses pasan por cada arco, no cuál bus. En cada nodo se resuelve una asignación entre arcos que entran y que salen. FIFO o LIFO bastarían con diésel, pero ignoran la energía. Por eso hay dos reglas según el nodo:
\begin{itemize}
  \item En un \textbf{electroterminal}: maximizar la espera mínima entre llegada y salida, para abrir ventanas de carga.
  \item En un \textbf{terminal convencional}: equilibrar el consumo acumulado entre rotaciones, para no juntar tramos que dejen un déficit irrecuperable.
\end{itemize}

\textbf{Fase 3: validación energética.} Recorre cada rotación actualizando el SOC según el consumo de cada tramo. Si hay un electroterminal cerca y la batería lo pide, inserta una carga en la holgura antes del siguiente viaje. Si el tramo previo es un traslado en vacío, lo adelanta para llegar antes y alargar la ventana. Respeta el mínimo de 10\,\%. Si aun así el SOC baja del mínimo, corta la rotación en el último punto desde donde puede ir a cargar, y \textbf{un bus nuevo toma el resto}.

\begin{minipage}[t]{0.48\linewidth}
\textbf{A favor}
\begin{itemize}
  \item Las líneas de tiempo quitan hasta el 97\,\% de los arcos vacíos y permiten calcular un límite inferior de vehículos. Eso importa porque la flota es el 66\,\% del costo.
  \item Las reglas de la Fase 2 ya piensan en eléctrico: abren ventanas de carga desde el inicio.
  \item Sectorizar por electroterminal sigue los cupos de cada uno, así que se pierde poco acoplamiento.
  \item Los buses que agrega la Fase 3 miden cuánta flota extra cuesta la congestión de los 700 cupos.
\end{itemize}
\end{minipage}\hfill
\begin{minipage}[t]{0.48\linewidth}
\textbf{En contra}
\begin{itemize}
  \item Supone carga lineal. Las baterías de litio cargan más lento cerca del 100\,\%.
  \item La asignación (Fase 1) se resuelve desacoplada de la batería (Fase 3): se pierde la optimalidad, y cortar rotaciones puede pedir más buses que un exacto.
  \item La escala comprobada es 10.710 viajes, así que dependemos de la sectorización previa, y la solución queda atada a ella.
\end{itemize}
\end{minipage}

\textbf{Versión de la presentación} (misma idea, otras cifras): la Fase 1 es un flujo de costo mínimo, polinomial, que resuelve 10.584 viajes en segundos; la Fase 2 produce del orden de 1.051 rotaciones; en la Fase 3, el 90,6\,\% de los viajes tiene holgura para una carga completa.

\subsection{Hallazgos de literatura de la presentación}
\begin{itemize}
  \item \textbf{Perumal, Lusby y Larsen (2022):} el E-VSP es una extensión reconocida del VSP. Los métodos van del flujo en redes a las metaheurísticas, el mismo origen de nuestras tres metodologías.
  \item \textbf{Alvo, Angulo y Klapp (2021):} electrificar sin capacidad de carga suficiente exige más buses. Con 29 buses y un solo cargador hacen falta 41 vehículos; con 3 cargadores bastan los 29. Hay que medir la congestión antes de decidir.
  \item \textbf{Othman et al. (2025):} unas rutas son más electrificables que otras. Para clasificarlas pesa más la velocidad operacional que la energía.
\end{itemize}

% ============================================================
\section{Decisiones y su justificación}
\label{sec:decisiones}

El anuncio de la prueba dice que hay que entender ``en detalle la justificación de las decisiones tomadas''. Esta tabla es esa lista: cada fila puede ser una pregunta.

\begingroup
\small
\begin{longtable}{@{}L{4.0cm}L{6.0cm}L{5.6cm}@{}}
\toprule
\hd{Decisión} & \hd{Por qué} & \hd{Qué sacrifica o arriesga} \\
\midrule
\endfirsthead
\toprule
\hd{Decisión} & \hd{Por qué} & \hd{Qué sacrifica o arriesga} \\
\midrule
\endhead
\bottomrule
\endfoot
Filtrar \arch{agency\_id = RM} y \arch{route\_type = 3} & El GTFS mezcla metro, trenes y aeropuerto; solo nos importan los buses RED & Nada relevante: se quitan 10 de 427 rutas \\
Expandir frecuencias con $t_n = \text{inicio} + n\cdot\text{headway}$ & \arch{exact\_times = 0}: no hay horarios exactos & Las salidas son representativas, no observadas \\
Consumo lineal de 1,4\,kWh/km & No se entregó consumo; 350\,kWh / 250\,km es lo único derivable & Ignora pendiente, clima, pasajeros, tráfico y degradación \\
SOC mínimo de 10\,\% & Límite de los datos, protege la vida útil de la batería & Deja 315\,kWh y unos 225\,km útiles \\
SOC inicial de 100\,\% & Los buses no están individualizados ni traen SOC & Optimista; la cadena semanal lo reemplaza por un estado que se hereda \\
Flota irrestricta & Aclaración docente; además 1.200 buses están muy bajo la cota de 6.601, así que como tope harían infactible el problema & La flota pasa a ser la variable que más pesa en el costo \\
Frecuencias exógenas & No hay datos de demanda ni de ocupación & No se puede ahorrar recortando servicio \\
Laboral, sábado y domingo por separado & Tres regímenes de demanda distintos (picos 08:00, 15:00 y 18:00) & Se encadenan con condiciones de borde para no perder el ciclo semanal \\
Recargas parciales y en cualquier electroterminal & Da flexibilidad y es técnicamente posible & Puede subestimar traslados si el electroterminal cercano está lleno \\
Capacidad eléctrica suficiente para cargar todo a la vez & No hay potencia de transformadores ni conteo validado de cargadores & Los 700 cupos físicos sí se miden, con el KPI de ocupación máxima \\
No modelar perturbaciones & No hay datos para caracterizar tráfico, accidentes ni fallas & Factible en el modelo no es lo mismo que factible en la calle \\
Cota de flota por simultaneidad & Se calcula directo de los datos, antes de tener modelo & Es floja: ignora traslados y recargas \\
KPIs sin metas y cobertura como factibilidad & Sin solución base, una meta sería arbitraria & Las metas quedan para la Entrega 2 \\
Red tiempo-espacio y no MILP ni columnas & Única con escala demostrada cerca de la nuestra & Sin garantía de óptimo \\
Sectorizar por electroterminal & Deja cada subproblema en unos 12.000 viajes, y el acoplamiento entre depósitos es débil & La solución depende de cómo se asignen las rutas a sectores \\
\end{longtable}
\endgroup

\subsection*{La cadena de cotas: la respuesta que muestra que entendiste}
Cada paso agrega una restricción, así que la flota solo puede subir:
\begin{enumerate}
  \item 6.601 buses: simultaneidad pura, sin traslados.
  \item Flujo sin batería sobre toda la red: agrega los traslados en vacío.
  \item Flujo sin batería por sector (nuestra Fase 1): agrega la sectorización.
  \item Óptimo real con batería y carga: agrega la energía.
  \item Nuestra solución: la heurística de las Fases 2 y 3 puede quedar sobre el óptimo.
\end{enumerate}
Por eso el Gap contra 6.601 sobrestima la distancia real al óptimo. La cota que da la Fase 1 es válida para el problema sectorizado, no para el global.

\subsection*{Por qué la presentación dijo ``polinomial'' y el informe dice ``MILP multibienes''}
Con un solo electroterminal por sector y un solo tipo de bus, el flujo es de un solo bien: un flujo de costo mínimo, cuya matriz es totalmente unimodular, así que el LP ya da una solución entera. Si en un sector hubiera varios depósitos o tipos de bus, vuelve a ser multibienes y entero. Las dos frases son coherentes; conviene decirlo así.

\begin{ojo}
La cadena de cotas, este argumento de unimodularidad y la mención a NP-dificultad de la sección~\ref{sec:kpis} son razonamientos para defender el informe, no están escritos en él. Úsalos como argumento propio, no como cita.
\end{ojo}

\subsection*{Cifras que no calzan entre documentos}
Si te las preguntan, usa la del informe y explica la diferencia sin dudar.
\begin{center}
\small
\begin{tabularx}{\linewidth}{@{}L{2.5cm}L{2.9cm}L{3.4cm}Y@{}}
\toprule
\hd{Dato} & \hd{Informe} & \hd{Presentación} & \hd{Cómo explicarlo} \\
\midrule
Corridas laborales & 64.502 & 64.603 (lámina de objetivo) & Versión anterior del procesamiento (418 rutas). Vale la del informe \\
Cota de flota & 6.601 & 6.601 en las láminas; 6.603 detrás del costo & El piso de costo usó $6.603 \times 250 = 1.650.750$ USD \\
Costo diario & No aparece & 2.512.204 (objetivo) y 2.545.946 (KPIs) & Dos pisos con supuestos de energía distintos. Lo robusto: unos 2,5 millones de USD al día, 66\,\% flota \\
de Vos et al.\ (2024) & 816 viajes en 7 horas & 816 viajes en 28 horas & Lo que importa es la escala: 79 veces más chica que la nuestra \\
Arcos del MILP & 4,09 millones de arcos viaje-viaje & 1.777 millones posibles, 23\,\% factibles & 23\,\% de 1.777 millones son unos 409 millones: probablemente ``4,09'' es 409 millones. Di ``cientos de millones de arcos factibles'' \\
\bottomrule
\end{tabularx}
\end{center}

% ============================================================
\section{Variantes del problema}
\label{sec:variantes}

La prueba puede pedir ``cómo aplicarían lo desarrollado a una variante''. La receta: di qué supuesto se rompe, en qué fase de nuestra metodología entra el cambio, y qué KPI lo mostraría.

\begingroup
\small
\begin{longtable}{@{}L{3.3cm}L{4.4cm}L{7.9cm}@{}}
\toprule
\hd{Variante} & \hd{Qué cambia en el problema} & \hd{Cómo lo adaptamos} \\
\midrule
\endfirsthead
\toprule
\hd{Variante} & \hd{Qué cambia en el problema} & \hd{Cómo lo adaptamos} \\
\midrule
\endhead
\bottomrule
\endfoot
La flota se limita a 1.200 buses & Con la cota de 6.601, cubrir el 100\,\% se vuelve infactible & Cambia el objetivo: maximizar cobertura, o penalizar corridas sin cubrir. En la Fase 1 se agrega un tope al flujo que sale de los depósitos, sumado entre sectores; eso vuelve a acoplar los 5 subproblemas \\
La capacidad eléctrica sí es limitada (potencia o cargadores) & No todos los buses pueden cargar a la vez: aparece congestión & En la Fase 3, antes de insertar una carga, revisar la ocupación del electroterminal en esa franja. Si está lleno: mover la carga a otra holgura, ir a otro electroterminal o agregar un bus. El KPI de ocupación máxima pasa a ser restricción. En columnas, es la restricción del maestro por electroterminal y bloque horario \\
Carga no lineal (más lenta cerca del 100\,\%) & Cargar de 80\,\% a 100\,\% cuesta mucho más tiempo que de 10\,\% a 30\,\% & En la Fase 3, calcular el tiempo con una curva lineal por tramos. Conviene preferir cargas parciales cortas en la zona rápida \\
Precio de la energía según la hora (0,10 a 0,22\,USD/kWh) & Importa cuándo se carga, no solo cuánto & En la Fase 3, entre las holguras factibles elegir la más barata, por ejemplo la noche. El KPI de costo energético lo mide. Ojo: nunca a costa de agregar un bus, que vale 250\,USD \\
Tiempos de viaje inciertos (congestión) & Un encadenamiento justo puede fallar en la calle & Exigir holguras mínimas entre corridas y reserva de SOC sobre el 10\,\%. Evaluar el plan con simulación Monte Carlo y medir atrasos y quiebres de cadena \\
Buses heterogéneos (distintas baterías) & Cada tipo tiene otra autonomía y costo & Un bien por tipo de bus en el flujo de la Fase 1, que pasa a ser multibienes de verdad. En la Fase 3, un consumo y un límite por tipo \\
Flota mixta diésel y eléctrica & El diésel no tiene restricción de energía & Asignar a diésel las cadenas que la Fase 3 no logra hacer factibles. Othman et al.\ (2025) sugiere electrificar primero las rutas más lentas \\
Un electroterminal fuera de servicio & Un sector se queda sin carga & Reasignar sus rutas a los sectores vecinos y volver a resolver. El KPI de ocupación muestra si los vecinos aguantan \\
SOC inicial distinto o desconocido & El 100\,\% ya no vale & Usar el estado heredado de la cadena semanal, o exigir que el día termine en el mismo estado en que empezó (el ciclo laboral) \\
Cambian las frecuencias & Otra instancia & Como las frecuencias son exógenas, se vuelve a expandir con $t_n$ y se resuelve de nuevo. El pipeline es reproducible \\
Instancia chica (un electroterminal, pocas corridas) & El tamaño deja de ser el problema & Resolver exacto con MILP o columnas y comparar con nuestra heurística: así se estima su gap real \\
\end{longtable}
\endgroup

\subsection*{Variante que vale preparar a fondo: capacidad de carga limitada}
Es la que más probablemente pregunten, porque el enunciado insiste en ella y nosotros la supusimos suficiente. La respuesta completa tiene tres partes:
\begin{enumerate}
  \item \textbf{La cuenta de por qué aprieta.}
  \[
    \underbrace{6.601 \times 0{,}91 \times 105\ \text{min}}_{\text{demanda de carga}} \approx 630.700\ \text{min}
    \qquad\text{vs.}\qquad
    \underbrace{700 \times 1.440\ \text{min}}_{\text{oferta de cupos}} = 1.008.000\ \text{min}.
  \]
  En promedio se usa el 63\,\%, pero la carga no puede repartirse pareja: se concentra fuera de las horas punta.
  \item \textbf{Dónde entra en el método:} la Fase 3, al insertar cada carga, revisa la ocupación por franja.
  \item \textbf{Qué muestra:} los buses que agrega la Fase 3 por falta de cupo miden el costo de la congestión. Es el mismo hallazgo de Alvo et al.\ (2021): 29 buses con un cargador necesitan 41 vehículos; con 3 cargadores bastan 29.
\end{enumerate}

% ============================================================
\section{Preguntas de práctica}
\label{sec:practica}

Diez preguntas del estilo que anunciaron. Escribe tu respuesta a mano en 12 minutos y después compara: ¿nombraste una cifra, justificaste contra una alternativa y dijiste qué se sacrifica? Para esconder las respuestas modelo, cambia \verb|\respuestastrue| por \verb|\respuestasfalse| en el preámbulo.

\pregunta{Describa el problema y explique qué lo hace distinto de programar una flota diésel.}
\respuesta{Hay que decidir cuántos buses usar, qué corridas hace cada uno y en qué orden, y cuándo y dónde recarga, para cubrir las 64.502 corridas laborales de la RED una vez cada una al mínimo costo. En diésel solo se decide la secuencia de viajes, porque cargar combustible es rápido y no tiene cupo. En eléctrico, un bus recorre en promedio 203,9\,km al día contra 225\,km útiles: casi una recarga diaria de 105 minutos, que solo puede hacerse en uno de 5 electroterminales con 700 cupos. Entonces la carga compite con las corridas por el tiempo del bus, y los buses compiten entre sí por los cupos. El plan de viajes y el de carga deben decidirse juntos.}

\pregunta{¿Cómo obtuvieron la cota de 6.601 buses? ¿Por qué no es la flota necesaria?}
\respuesta{Se contó cuántas corridas reconstruidas están activas al mismo tiempo en cada minuto del día laboral. El máximo es 6.601, a las 08:00. Como un bus no puede hacer dos corridas a la vez, ningún plan puede usar menos: es condición necesaria. No es suficiente porque ignora los traslados en vacío entre el fin de una corrida y el inicio de otra, las recargas de 105 minutos y la regulación. Todo eso saca buses de servicio, así que la flota factible será mayor. Por eso el Gap contra 6.601 no es un gap de optimalidad.}

\pregunta{Explique cómo reconstruyeron las corridas a partir del GTFS y qué limitación introduce.}
\respuesta{Primero se filtró \arch{agency\_id = RM} y \arch{route\_type = 3}, quedando 417 recorridos y 6.620 patrones laborales. Como todas las frecuencias RED tienen \arch{exact\_times = 0}, el GTFS no da horarios sino reglas: una ventana y un headway. Se generaron salidas con $t_n = \text{inicio} + n\cdot\text{headway}$ mientras $t_n$ sea menor que el fin de la ventana. Por ejemplo, de 07:00 a 08:00 con 900\,s salen 07:00, 07:15, 07:30 y 07:45. Eso dio 64.502 corridas. La duración sale de la primera y última parada, y la distancia de \arch{shape\_distances}. La limitación: son horarios representativos de la frecuencia declarada, no observados, y las duraciones son nominales, sin congestión.}

\pregunta{¿Por qué descartaron el MILP monolítico y la generación de columnas?}
\respuesta{El MILP es el único que garantiza el óptimo y es simple de escribir, pero al expandir los arcos entre corridas por niveles de SOC y por electroterminal llega a unos $5\times10^{8}$ variables binarias, fuera del alcance de un solver. Además, si no termina no entrega nada aprovechable. La generación de columnas es mejor: da una cota inferior de calidad, porque cada columna ya es una jornada factible. Pero su escala publicada es de 816 viajes (de Vos et al., 2024), 79 veces menos que nuestras 64.502 corridas. El cuello de botella es el \emph{pricing}, un camino más corto con recursos sobre una red expandida por SOC, que corre cientos o miles de veces.}

\pregunta{Explique las tres fases de la metodología escogida y qué se relaja en cada una.}
\respuesta{Fase 1: un flujo sobre una red tiempo-espacio agregada con líneas de tiempo, que minimiza el costo de vehículos con cobertura exacta. Relaja la batería y la capacidad de carga. Entrega flota, km en vacío y flujo por arco. Fase 2: descompone el flujo en rotaciones, es decir jornadas de un bus concreto, resolviendo una asignación en cada nodo. En electroterminales maximiza la espera mínima para abrir ventanas de carga; en terminales convencionales equilibra el consumo acumulado. Fase 3: recorre cada rotación actualizando el SOC e inserta cargas en las holguras, adelantando traslados en vacío si hace falta. Si igual cae bajo 10\,\%, corta la rotación y agrega un bus. Lo que se relaja en la 1 se recupera en la 3, a costa de la optimalidad.}

\pregunta{¿Por qué sectorizar por electroterminal? ¿Qué se pierde?}
\respuesta{Olsen et al.\ (2020) demuestran la metodología hasta 10.710 viajes; nuestras 64.502 corridas son 6 veces más. Repartir las 417 rutas entre los 5 electroterminales por proximidad deja subproblemas de unos 12.000 viajes, dentro de la escala comprobada, y se pueden resolver por separado. Es natural porque cada electroterminal tiene sus propios cupos, y en el MDVSP clásico el acoplamiento entre depósitos es débil cuando cada uno tiene su capacidad (Kliewer et al., 2006). Se pierde que un bus no puede cruzar sectores: un traslado o una carga más barata en el electroterminal vecino ya no está disponible. La solución queda atada a esa asignación y la cota de la Fase 1 vale para el problema sectorizado, no para el global.}

\pregunta{Elija dos KPIs, justifíquelos y explique por qué la cobertura no es un KPI.}
\respuesta{Número de buses: la flota es el 66\,\% del piso de costo (1.650.750 de 2.512.204\,USD/día), así que es el indicador que más mueve el resultado. Se contextualiza con $\text{Gap}_F = (F_{sol} - 6.601)/6.601$. Porcentaje de km sin pasajeros: los traslados en vacío son el costo que la cota de simultaneidad no ve, y solo el 3,1\,\% de los encadenamientos posibles es barato. Este KPI muestra si el plan es eficiente o si arrastra buses por la ciudad. La cobertura del 100\,\% no es KPI porque no admite grados: si falta una corrida, el plan es infactible, no peor. Lo mismo la factibilidad temporal, espacial y energética: son condiciones de entrada, no desempeño.}

\pregunta{Si el equipo docente impusiera un máximo de 1.200 buses, ¿qué cambia?}
\respuesta{El problema se vuelve infactible tal como está: la cota por simultaneidad es 6.601, cinco veces más. Habría que cambiar el objetivo: maximizar corridas cubiertas, o minimizar costo más una penalidad por corrida sin cubrir. Ahí la cobertura sí pasa a ser un KPI. En nuestra metodología, la Fase 1 lleva un tope al flujo total que sale de los depósitos. Ese tope suma los 5 sectores, así que vuelve a acoplarlos. Se podría repartir el tope entre sectores en proporción a su demanda en hora punta, o iterar el reparto. Por esto también tiene sentido el supuesto de flota irrestricta.}

\pregunta{Si la potencia de los electroterminales no alcanzara para cargar a todos a la vez, ¿cómo lo incorpora?}
\respuesta{Nuestro supuesto actual es capacidad eléctrica suficiente, porque no hay datos de transformadores. Si se levanta, la Fase 3 debe llevar la ocupación de cada electroterminal por franja horaria. Al insertar una carga, revisa si hay cupo y potencia. Si no, mueve la carga a otra holgura, la manda a otro electroterminal o, si nada funciona, corta la rotación y agrega un bus. Los buses agregados miden el costo de la congestión, como en Alvo et al.\ (2021): 29 buses con un cargador necesitan 41 vehículos. También conviene que la Fase 2 ordene rotaciones para escalonar las llegadas al electroterminal. El KPI de ocupación máxima pasa a ser restricción.}

\pregunta{Si su metodología no garantiza optimalidad, ¿cómo sabrán si la solución es buena?}
\respuesta{Con tres referencias. Primero, la cota inferior: la flota nunca puede bajar de 6.601, y la Fase 1 da otra más ajustada para el problema sectorizado. El Gap contra ellas acota la distancia al óptimo por arriba. Segundo, una solución base sencilla y factible, contra la cual se comparan todos los KPIs. Tercero, en una instancia chica, por ejemplo un electroterminal y pocas corridas, se resuelve exacto con MILP o columnas y se compara con la heurística: eso estima su gap real. Además, el análisis de sensibilidad sobre consumo, SOC inicial y capacidad de carga muestra si las conclusiones dependen de supuestos frágiles.}

% ============================================================
\section{Repaso final}
\label{sec:repaso}

Para leer mañana antes de entrar: las cifras que conviene poder escribir sin mirar y seis frases que resumen el informe.

\begingroup
\small
\begin{longtable}{@{}L{3.4cm}L{7.0cm}L{5.2cm}@{}}
\toprule
\hd{Cifra} & \hd{Qué es} & \hd{De dónde sale} \\
\midrule
\endfirsthead
\toprule
\hd{Cifra} & \hd{Qué es} & \hd{De dónde sale} \\
\midrule
\endhead
\bottomrule
\endfoot
417 de 427 & Recorridos RED & Filtro \arch{agency\_id = RM}, \arch{route\_type = 3} \\
64.502 & Corridas de un día laboral (43.619 sábado, 37.377 domingo) & Expansión de frecuencias \\
6.601 & Cota inferior de flota laboral, pico a las 08:00 & Máximo de corridas simultáneas \\
1.345.413\,km y 1.884\,MWh & Km de servicio y energía de un día laboral & Suma por corrida $\times$ 1,4\,kWh/km \\
20,9\,km, 78,6\,min, 29,2\,kWh & Corrida laboral promedio & Tabla 1.4 del informe \\
350\,kWh, 250\,km, 180\,kW & Batería, autonomía nominal y potencia de carga & \arch{vehicles.csv} \\
1,4\,kWh/km & Consumo supuesto lineal & 350 / 250 \\
10\,\%, 315\,kWh, 225\,km & SOC mínimo, energía y autonomía útiles & 90\,\% de 350\,kWh \\
105\,min & Una recarga completa & 315\,kWh / 180\,kW \\
203,9\,km y 91\,\% & Km por bus al día y fracción de la autonomía & 1.345.413\,km / unos 6.600 buses \\
5 y 700 & Electroterminales y cupos físicos & \arch{depots.csv} \\
1.200 & Buses de referencia, no es tope & \arch{parameters.csv} y aclaración docente \\
250 $\cdot$ 0,5 $\cdot$ 0,03 $\cdot$ 5 & USD por bus, por km, por minuto de espera y por carga & \arch{parameters.csv} \\
0,10 a 0,22 & USD por kWh según la hora & \arch{electricity\_prices.csv} \\
2.512.204 y 66\,\% & Piso de costo diario y peso de la flota en él & Presentación, lámina de objetivo \\
1.777 millones, 23\,\%, 3,1\,\% & Encadenamientos posibles, geográficamente factibles y baratos & Presentación, lámina de dificultad \\
$5\times10^{8}$ & Variables binarias del MILP completo & Informe, sección 2.1 \\
816 y 79 veces & Escala de columnas (de Vos) y cuánto más grande es nuestra instancia & Informe, sección 2.2 \\
10.710 y 97\,\% & Escala de Olsen y arcos vacíos que eliminan las líneas de tiempo & Informe, sección 2.3 \\
Unos 12.000 & Viajes por sector tras sectorizar & 64.502 corridas / 5 electroterminales \\
\end{longtable}
\endgroup

\begin{clave}
\textbf{Seis frases}
\begin{enumerate}
  \item El plan de viajes y el plan de carga se deciden juntos o no se deciden.
  \item 6.601 es un piso: necesario, no suficiente.
  \item La flota es el 66\,\% del costo, así que el objetivo primario es usar menos buses, no gastar menos energía.
  \item La cobertura del 100\,\% es condición de factibilidad, no un KPI.
  \item Sacrificamos la garantía de optimalidad a cambio de una escala demostrada.
  \item Hay que distinguir factibilidad dentro del modelo de factibilidad operacional real.
\end{enumerate}
\end{clave}

\subsection*{Las cuatro referencias del informe}
Por si piden citar:
\begin{itemize}
  \item Alvo, M., Angulo, G. y Klapp, M. (2021). \emph{An exact solution approach for an electric bus dispatch problem}. Escuela de Ingeniería, Pontificia Universidad Católica de Chile. El exacto de la UC.
  \item de Vos, M.\,H., van Lieshout, R.\,N. y Dollevoet, T. (2024). \emph{Electric vehicle scheduling in public transit with capacitated charging stations}. Transportation Science. Columnas con estaciones de carga capacitadas.
  \item Kliewer, N., Mellouli, T. y Suhl, L. (2006). \emph{A time--space network based exact optimization model for multi-depot bus scheduling}. European Journal of Operational Research, 175(3), 1616--1627. Red tiempo-espacio multidepósito.
  \item Olsen, N., Kliewer, N. y Wolbeck, L. (2020). \emph{A study on flow decomposition methods for scheduling of electric buses in public transport based on aggregated time-space network models}. Central European Journal of Operations Research. La base de nuestra metodología.
\end{itemize}

\end{document}
